\documentclass[12pt,a4paper]{article}

% Biên dịch bằng XeLaTeX:
%   latexmk -xelatex Ke_hoach_nghien_cuu_HCORAP.tex

\usepackage{fontspec}
\setmainfont{Times New Roman}
\setsansfont{Arial}
\setmonofont{Courier New}

\usepackage[a4paper,margin=2.5cm]{geometry}
\usepackage{amsmath,amssymb,mathtools}
\usepackage{booktabs,longtable,tabularx,array,multirow}
\usepackage{enumitem}
\usepackage{xcolor}
\usepackage{microtype}
\usepackage{fancyhdr}
\usepackage{float}
\usepackage{hyperref}

\hypersetup{
  colorlinks=true,
  linkcolor=blue!55!black,
  citecolor=blue!55!black,
  urlcolor=blue!65!black,
  pdftitle={Đánh giá đa mục tiêu có kiểm soát cho HCORAP bằng MaxSAT},
  pdfauthor={Người thực hiện}
}

\pagestyle{fancy}
\fancyhf{}
\lhead{Kế hoạch nghiên cứu HCORAP}
\rhead{Weighted/Lex/$\varepsilon$-constraint MaxSAT}
\cfoot{\thepage}
\setlength{\headheight}{16pt}

\setlist[itemize]{leftmargin=1.6em,itemsep=0.3em,topsep=0.4em}
\setlist[enumerate]{leftmargin=1.8em,itemsep=0.35em,topsep=0.4em}
\renewcommand{\arraystretch}{1.2}

\newcolumntype{Y}{>{\raggedright\arraybackslash}X}
\newcommand{\SIM}{\mathrm{SIM}}
\newcommand{\STAB}{\mathrm{STAB}}
\newcommand{\CONT}{\mathrm{CONT}}
\newcommand{\OT}{\mathrm{OT}}
\newcommand{\COV}{\mathrm{COV}}
\newcommand{\SEQ}{\mathrm{SEQ}}
\newcommand{\HN}{\mathrm{HN}}
\newcommand{\HE}{\mathrm{HE}}

\title{%
  \textbf{KẾ HOẠCH NGHIÊN CỨU VÀ TRIỂN KHAI}\\[0.6em]
  \Large Đánh giá đa mục tiêu có kiểm soát cho bài toán phân bổ nguồn lực
  trong dịch vụ chăm sóc tại nhà (HCORAP)
}
\author{Người thực hiện: \textit{[Điền họ và tên]}\\
        Đơn vị: \textit{[Điền đơn vị công tác/học tập]}}
\date{\today}

\begin{document}

\maketitle
\thispagestyle{empty}

\begin{abstract}
Tài liệu này trình bày kế hoạch nghiên cứu nhằm đánh giá và cải tiến mô hình
Home Care Optimal Resource Allocation Problem (HCORAP) được giải bằng
Weighted Partial MaxSAT. Nghiên cứu không mặc định một phương pháp tốt hơn
phương pháp khác. Thay vào đó, baseline weighted đã hiệu chỉnh, hai chính sách
ưu tiên theo thứ tự từ điển và phương pháp
$\varepsilon$-constraint được so sánh trên benchmark đã kiểm toán và hiệu chỉnh.
Thiết kế này tách riêng ảnh hưởng của lỗi benchmark, cách đặt trọng số và phương
pháp tối ưu. Nghiên cứu giữ nguyên phạm vi assignment--scheduling theo tuần,
không mở rộng sang routing hoặc uncertainty trong giai đoạn chính. Kết quả dự
kiến gồm baseline có thể tái lập, benchmark HCORAP phiên bản 2, các bộ giải
weighted/lexicographic/$\varepsilon$-constraint MaxSAT, ít nhất một baseline
ngoài MaxSAT, bộ kiểm chứng nghiệm độc lập và đánh giá thực nghiệm có kiểm soát.
Ba trục cải tiến encoding gồm cardinality encoding, implied constraints và
symmetry breaking được cấu hình độc lập để đo ảnh hưởng đến kích thước công
thức và thời gian chứng minh tối ưu mà không trộn với ảnh hưởng của policy tối ưu.
\end{abstract}

\vfill
\noindent\textbf{Từ khóa:} HCORAP; MaxSAT; lexicographic optimization;
$\varepsilon$-constraint; multi-objective optimization; home care scheduling.

\clearpage
\tableofcontents
\clearpage

\section{Bối cảnh và động cơ nghiên cứu}

Nghiên cứu gốc mô hình hóa việc phân công caregiver cho các dịch vụ chăm sóc
tại nhà bằng Weighted Partial MaxSAT. Một nghiệm phải thỏa tất cả hard
constraints về coverage, availability, qualification, simultaneity và giới hạn
giờ làm; đồng thời tối đa hóa ba thành phần:
\begin{equation}
  \text{similarity}+\text{stability}-\text{cost}.
  \label{eq:original-objective}
\end{equation}

Mã C++ chính thức được xem là baseline của tác giả. Để phép đo hiệu năng không
bị nhiễu bởi khác biệt ngôn ngữ, weighted baseline tương đương, hai chính sách
lexicographic và phương pháp $\varepsilon$-constraint đều được triển khai bằng
C++ và gọi cùng một binary MaxSAT. Các tệp Python tự cài đặt chỉ đóng vai trò
oracle, generator, verifier thứ hai và test harness; runtime Python không được
đưa vào bảng so sánh hiệu năng chính.

Việc kiểm toán bài báo, generator và 800 instance chính thức cho thấy các vấn
đề đáng chú ý sau:
\begin{itemize}
  \item tham số $V$ xác định tổng $S=U\times V$, nhưng generator không bảo đảm
        mỗi user có đúng $V$ services;
  \item thuộc tính ngôn ngữ của user được sinh từ tập nhãn race, làm cho thành
        phần language compatibility không hoạt động đúng;
  \item một số instance có service không tồn tại bất kỳ cặp agent--time-slot
        khả thi nào, dẫn đến trivial UNSAT;
  \item các cấu hình có số agent khác nhau được sinh độc lập, vì vậy chưa tạo
        thành thí nghiệm paired/nested để suy luận tác động của $A$;
  \item similarity và stability được báo cáo bằng tổng thô, nên tăng theo kích
        thước instance và không phù hợp để so sánh giữa các giá trị $S$;
  \item cost gần bằng 0 trên phần lớn cấu hình, nên benchmark chưa tạo đủ
        trade-off giữa similarity, continuity và overtime;
  \item nghiên cứu chưa đánh giá độ nhạy theo trọng số và chưa so sánh với một
        cách tối ưu đa mục tiêu không cần cộng các tiêu chí bằng weighted sum.
  \item sorting network được dùng cố định cho cardinality constraints, nên chưa
        biết Totalizer hoặc các strengthening bảo toàn optimum có cải thiện
        propagation và time-to-proof hay không;
  \item detector exact trên 800 instance tìm thấy time-slot symmetry có candidate
        thực sự trong 553 instance ($69.1\%$), service symmetry trong 340 instance
        ($42.5\%$), nhưng không tìm thấy agent symmetry hoàn toàn.
\end{itemize}

Những quan sát trên dẫn tới nhu cầu xây dựng lại quy trình đánh giá trước khi
mở rộng mô hình sang các yếu tố phức tạp hơn như routing, uncertainty hoặc
rescheduling.

\section{Hướng nghiên cứu được lựa chọn}

Hướng nghiên cứu chính được xác định như sau:

\begin{quote}
\textbf{Đánh giá có kiểm soát calibrated Weighted MaxSAT, sequential
Lexicographic MaxSAT và similarity-budget $\varepsilon$-constraint MaxSAT cho
HCORAP trên benchmark đã hiệu chỉnh, nhằm xác định khi nào mỗi cách biểu diễn
ưu tiên tạo ra phương án phân công phù hợp.}
\end{quote}

Đây là hướng khả thi vì:
\begin{enumerate}
  \item giữ nguyên cấu trúc biến và phần lớn hard constraints của mô hình gốc;
  \item tận dụng được C++ gốc để tái lập baseline và mở rộng trực tiếp cùng
        C++ model cho các phương pháp mới;
  \item không phụ thuộc ngay vào dữ liệu GPS hoặc dữ liệu vận hành thời gian thực;
  \item cho phép đánh giá cải tiến bằng các KPI định lượng rõ ràng;
  \item vẫn giữ MaxSAT là đối tượng nghiên cứu trung tâm, đồng thời có baseline
        ngoài MaxSAT để tránh kết luận chỉ dựa trên một hệ giải.
\end{enumerate}

\subsection{Định vị đóng góp và tuyên bố tính mới}

Tối ưu theo thứ tự từ điển không phải ý tưởng mới trong scheduling hoặc home
care. Weighted Partial MaxSAT cũng đã được áp dụng cho staff scheduling
\cite{demirovic2019,mosquera2019,bodvarsdottir2025}. Multi-objective MaxSAT hiện
đã có các thuật toán và solver chuyên biệt \cite{jabs2025}. Vì vậy, nghiên cứu
\emph{không} tuyên bố ``lần đầu dùng lexicographic optimization cho scheduling''
hoặc ``lần đầu dùng tối ưu đa mục tiêu cho home care''.

Đóng góp được định vị ở tổ hợp có thể kiểm chứng sau:
\begin{enumerate}
  \item kiểm toán và tái lập formulation HCORAP/MaxSAT gốc;
  \item xây dựng benchmark paired/nested đã sửa lỗi và có mức tải được hiệu
        chỉnh trên calibration seeds độc lập;
  \item so sánh calibrated weighted sum với các chính sách ưu tiên tuyệt đối và
        $\varepsilon$-constraint trên cùng Boolean model;
  \item đánh giá riêng Totalizer, implied constraints và symmetry breaking bằng
        ma trận ablation có baseline \texttt{none} rõ ràng;
  \item xác định phần cải thiện đến từ benchmark, metric hay optimizer bằng
        ablation và ít nhất một baseline CP-SAT hoặc MILP;
  \item công bố verifier, cấu hình, raw results và các kết quả âm nếu calibrated
        weighted sum đã đủ tốt.
\end{enumerate}

Việc tiếp tục dùng MaxSAT cho phép giữ nguyên phần lớn encoding gốc và thực hiện
so sánh có kiểm soát. Nó không phải bằng chứng MaxSAT tốt hơn CP-SAT/MILP; câu
hỏi này phải được trả lời bằng thực nghiệm. Nếu calibrated weighted baseline
cho chất lượng và độ ổn định tương đương các phương pháp còn lại, kết luận đúng
là hạn chế chủ yếu nằm ở benchmark, không phải ở weighted sum.

\subsection{Phạm vi}

Nghiên cứu giữ nguyên các giả định sau:
\begin{itemize}
  \item lập lịch theo tuần;
  \item mỗi service chiếm một time slot;
  \item mỗi service được thực hiện bởi tối đa một agent tại một time slot;
  \item qualification, agent availability và service availability là dữ liệu
        xác định;
  \item continuity được định nghĩa trên các sequence $\SEQ(q)$;
  \item số service được giao cho agent không vượt $\HN(a)+\HE(a)$.
\end{itemize}

Trong giai đoạn này, nghiên cứu \emph{không} bổ sung routing, travel time,
uncertainty, dynamic rescheduling, multi-sector coordination hoặc fairness
objective mới. Đây là các hướng tiếp theo sau khi baseline và objective hiện
tại đã được đánh giá đáng tin cậy.

\section{Câu hỏi và giả thuyết nghiên cứu}

\subsection{Câu hỏi nghiên cứu}

\begin{description}[leftmargin=3.4em,labelwidth=2.8em]
  \item[RQ1.] Nghiệm của weighted MaxSAT hiện tại nhạy như thế nào đối với
              trọng số continuity và overtime?
  \item[RQ2.] Với mức suy giảm similarity bị giới hạn trong $1$--$5\%$, có thể
              cải thiện continuity và overtime ở mức nào?
  \item[RQ3.] Sequential lexicographic và $\varepsilon$-constraint MaxSAT có
              cung cấp bảo đảm ưu tiên hoặc nghiệm không bị trội với chi phí
              tính toán nào so với calibrated Weighted MaxSAT và baseline
              ngoài MaxSAT?
  \item[RQ4.] Các kết luận của bài báo về similarity, stability và cost có còn
              giữ nguyên khi dùng benchmark paired/nested và metric chuẩn hóa
              hay không?
  \item[RQ5.] Totalizer, implied constraints và symmetry breaking ảnh hưởng như
              thế nào đến số biến, số clause, thời gian encode, time-to-proof và
              PAR-2 khi objective, solver và policy được giữ cố định?
\end{description}

\subsection{Giả thuyết}

\begin{description}[leftmargin=3.4em,labelwidth=2.8em]
  \item[H1.] Trên critical-load instances, thay đổi trọng số
             continuity/overtime làm thay đổi đáng kể cấu trúc nghiệm, dù tổng
             weighted objective có thể thay đổi ít; trên low-load instances,
             ảnh hưởng có thể không đáng kể.
  \item[H2.] Tồn tại các nghiệm giảm continuity penalty với tổn thất similarity
             nhỏ, đặc biệt trên các instance có nhiều lựa chọn agent.
  \item[H3.] Trade-off overtime chỉ thể hiện rõ trên các instance critical-load;
             benchmark có capacity dư sẽ không đủ để kiểm nghiệm cost objective.
  \item[H4.] Similarity trên mỗi service phụ thuộc chủ yếu vào số agent khả dụng,
             thay vì tăng đơn thuần theo tổng số services.
  \item[H5.] Không có một strengthening tốt nhất cho mọi instance; lợi ích của
             slot/service symmetry và implied constraints phụ thuộc candidate
             density, load và cấu trúc \texttt{SU/SEQ}, nên phải được chọn bằng paired
             pilot thay vì bật mặc định.
\end{description}

\section{Định nghĩa lại các tiêu chí đánh giá}

\subsection{Similarity}

Giữ nguyên tổng similarity:
\begin{equation}
  \SIM=\sum_{a=1}^{A}\sum_{s=1}^{S}r(a,s)y_{a,s}.
  \label{eq:sim}
\end{equation}

Để so sánh giữa các kích thước instance, sử dụng chỉ số chuẩn hóa:
\begin{equation}
  \SIM_N=\frac{\SIM}{r_{\max}S},
  \qquad r_{\max}=4.
  \label{eq:sim-normalized}
\end{equation}

Do $r_{\max}S$ có thể không đạt được khi qualification và time windows hạn
chế, báo cáo thêm chỉ số theo upper bound khả đạt:
\begin{equation}
  U_{\SIM}=\sum_{s=1}^{S}\max_{a\in\mathcal A_s}r(a,s),
  \qquad
  \SIM_A=\frac{\SIM}{U_{\SIM}},
  \label{eq:sim-attainable}
\end{equation}
trong đó $\mathcal A_s$ là tập agent có ít nhất một time slot khả thi cho
service $s$. Chỉ số này chỉ dùng khi $U_{\SIM}>0$. Trên các instance dùng similarity budget, báo cáo thêm
$\SIM/\SIM^*$ để đo trực tiếp mức mất mát so với similarity tối ưu.

\subsection{Continuity}

Đặt $D_q$ là số agent khác nhau tham gia sequence $q$:
\begin{equation}
  D_q=\sum_{a=1}^{A}ss_{a,q}.
  \label{eq:distinct-agents}
\end{equation}

Thay vì báo cáo stability reward có chứa hằng số, định nghĩa continuity penalty:
\begin{equation}
  \CONT=\sum_{q:\,|\SEQ(q)|>1}(D_q-1).
  \label{eq:continuity-penalty}
\end{equation}

\paragraph{Mệnh đề 1 (tương đương với stability gốc).}
Giả sử coverage là hard constraint và mọi sequence không rỗng. Đặt
$n_q=|\SEQ(q)|$ và
\begin{equation}
  \STAB=\sum_q\sum_{i=1}^{n_q}(1-c_{q,i}),
  \qquad
  C_{\mathrm{inst}}=\sum_q(n_q-1).
\end{equation}
Khi đó:
\begin{equation}
  \STAB=C_{\mathrm{inst}}-\CONT.
  \label{eq:stab-cont-equivalence}
\end{equation}

\emph{Chứng minh.} Theo định nghĩa của encoding gốc,
$c_{q,i}=1$ khi và chỉ khi $D_q\geq i$. Vì vậy, trong $n_q$ biến đếm có đúng
$D_q$ biến nhận giá trị 1 và
$\sum_{i=1}^{n_q}(1-c_{q,i})=n_q-D_q$. Mặt khác, sequence đơn có
$D_q=1$ nên đóng góp 0; do đó
$\CONT=\sum_q(D_q-1)$. Suy ra
$\STAB=\sum_q(n_q-D_q)=\sum_q(n_q-1)-\CONT$. \hfill$\square$

Hằng số đúng là $\sum_q(n_q-1)$, không phải $\sum_q n_q$. Tính tương đương
này chỉ được sử dụng khi mọi service đều được phục vụ.

Chỉ số continuity chuẩn hóa là:
\begin{equation}
  \CONT_N=
  \begin{cases}
    1-\dfrac{\sum_{q:\,|\SEQ(q)|>1}(D_q-1)}
                 {\sum_{q:\,|\SEQ(q)|>1}(|\SEQ(q)|-1)},
       & \text{nếu mẫu số}>0,\\[1.1em]
    1, & \text{nếu mẫu số}=0.
  \end{cases}
  \label{eq:continuity-normalized}
\end{equation}

Giá trị $\CONT_N=1$ thể hiện continuity tốt nhất: mỗi sequence chỉ sử dụng một
caregiver.

Khi coverage được nới thành soft objective trong stress test overload, có thể
có $D_q=0$ và Eq.~\eqref{eq:continuity-penalty} không còn thích hợp. Đặt
\begin{equation}
  K_q=\sum_{s\in\SEQ(q)}\mathrm{performed}(s),
  \qquad
  \CONT^{\mathrm{serv}}=\sum_q\max(0,D_q-1).
\end{equation}
Chỉ số continuity có điều kiện trên các service đã phục vụ là:
\begin{equation}
  \CONT_N^{\mathrm{serv}}=
  \begin{cases}
    1-\dfrac{\sum_q\max(0,D_q-1)}
                   {\sum_q\max(0,K_q-1)},
      & \text{nếu mẫu số}>0,\\[1.1em]
    1, & \text{nếu mẫu số}=0.
  \end{cases}
  \label{eq:continuity-served}
\end{equation}
Metric này chỉ được so sánh giữa các nghiệm có cùng $\COV$; đồng thời phải báo
cáo số sequence không được phục vụ và số sequence chỉ được phục vụ một phần.
Overload là stress test thứ cấp, không phải tập chính để kết luận về continuity.

\subsection{Overtime}

Số service vượt giờ chuẩn của agent $a$ là:
\begin{equation}
  e_a=\max\left(0,\sum_{s=1}^{S}y_{a,s}-\HN(a)\right).
  \label{eq:agent-overtime}
\end{equation}

Tổng overtime:
\begin{equation}
  \OT=\sum_{a=1}^{A}e_a,
  \qquad
  \OT_N=\frac{\OT}{S}.
  \label{eq:overtime}
\end{equation}

Ngoài $\OT_N$, cần báo cáo số agent có overtime, overtime cực đại trên một
agent và phân phối workload giữa các agent.

\subsection{Coverage}

Với benchmark low/critical, coverage tiếp tục là hard constraint và không cần
đưa $\COV$ vào objective. Chỉ trong stress test overload mới nới coverage thành
soft objective và định nghĩa:
\begin{equation}
  \COV=\sum_{s=1}^{S}\mathrm{performed}(s),
  \qquad
  \COV_N=\frac{\COV}{S}.
  \label{eq:coverage}
\end{equation}

Coverage phải là mục tiêu ưu tiên cao nhất khi cho phép bỏ một số services.
Mọi so sánh về $\SIM$, $\CONT$ và $\OT$ trong stress test phải condition trên
cùng mức $\COV$.

\section{Các phương pháp cần triển khai}

\subsection{B0: Weighted MaxSAT gốc}

Đặt $p=|P|$ là chi phí dương của một giờ vượt giờ. Baseline được viết dưới
dạng penalty tương đương (bỏ hằng số theo
Eq.~\eqref{eq:stab-cont-equivalence}):
\begin{equation}
  \max\; \SIM-w_c\CONT-w_o p\OT.
  \label{eq:weighted-baseline}
\end{equation}

Trong WCNF, similarity clauses $\langle y_{a,s},r(a,s)\rangle$ giữ nguyên;
stability clauses dùng $\langle\neg c_{q,i},w_c\rangle$; overtime clauses dùng
$\langle\neg w_{a,i},w_o p\rangle$. Không scalarize ở ngoài solver. Tệp cấu
hình phải lưu cả $(w_c,w_o,p)$ và các integer weights thực tế đã ghi vào WCNF.

Grid screening ban đầu là:
\begin{equation}
  w_c\in\{1,2,4,8\},
  \qquad
  w_o\in\{1,2,4,8\}.
\end{equation}

Đây là log-grid thăm dò, không phải khẳng định 16 cặp đã bao phủ đầy đủ không
gian trọng số. Sau screening, chỉ tinh chỉnh các khoảng mà assignment hoặc
metric thay đổi. Cấu hình gốc $(w_c,w_o)=(1,1)$ phải chạy trên toàn bộ official
benchmark; grid đầy đủ chỉ chạy trên subset phân tầng để giữ computational
budget khả thi. Mục đích của B0 là xây dựng calibrated weighted baseline mạnh,
không giả định $(1,1)$ là lựa chọn đúng duy nhất.

Calibrated weighted sum và lexicographic optimization giải quyết hai loại yêu
cầu khác nhau. Weighted sum cho phép bù trừ giữa các tiêu chí. Lexicographic
optimization bảo đảm không một tổng cải thiện ở tầng thấp nào đánh đổi một đơn
vị ở tầng cao. Weighted sum có thể mô phỏng thứ tự lexicographic bằng các
domination weights lớn hơn toàn bộ miền giá trị của các tầng dưới, nhưng các
trọng số đó phụ thuộc kích thước instance và khi ấy không còn biểu diễn trade-off
bù trừ thông thường. Vì vậy B1/B2 là các chính sách so sánh, không được mặc định
là ``tốt hơn'' B0.

\subsection{B1: Lexicographic MaxSAT}

Lexicographic optimization được thực hiện theo nhiều tầng, không dùng một
``big weight'' duy nhất. Để tránh gán nhãn giá trị khi chưa có stakeholder
study, hai kịch bản được gọi trung tính là:

\begin{align}
  \text{Continuity-priority:}\quad
    &-\CONT \succ \SIM \succ -\OT,
    \label{eq:continuity-priority}\\
  \text{Overtime-priority:}\quad
    &-\OT \succ -\CONT \succ \SIM.
    \label{eq:overtime-priority}
\end{align}

Ký hiệu $\succ$ nghĩa là mục tiêu bên trái có ưu tiên tuyệt đối cao hơn. Quy
trình tổng quát trên full-coverage instances:
\begin{enumerate}
  \item tối ưu mục tiêu ở tầng hiện tại;
  \item ghi nhận giá trị tối ưu;
  \item bổ sung bound bảo toàn giá trị tối ưu đó;
  \item chuyển sang mục tiêu tiếp theo.
\end{enumerate}

Ví dụ cho chính sách overtime-priority:
\begin{align*}
  &\min \OT  &&\Rightarrow \OT^*,\\
  &\min \CONT&&\text{s.t. }\OT\leq\OT^* \Rightarrow \CONT^*,\\
  &\max \SIM &&\text{s.t. }\OT\leq\OT^*,\ \CONT\leq\CONT^*.
\end{align*}

Vì $\OT^*$ và $\CONT^*$ đã tối ưu, các bất đẳng thức trên tương đương logic
với equality nhưng thường cần encoding đơn giản hơn. Stress test overload thêm
tầng $\max\COV$ đầu tiên và bound $\COV\geq\COV^*$ cho mọi tầng sau. Kiểm thử
phải bao gồm nhiều nghiệm đồng tối ưu để xác nhận tầng sau chọn đúng tie mà
không làm xấu tầng trước.

\subsection{B2: Similarity-budget
\texorpdfstring{$\varepsilon$}{epsilon}-constraint MaxSAT}

Đầu tiên tính similarity tối đa:
\begin{equation}
  \SIM^*=\max \SIM.
\end{equation}
Phân tích chính dùng full-coverage instances. Nếu B2 được dùng trong overload
stress test, phải tối đa hóa và cố định $\COV^*$ trước, rồi mới tính $\SIM^*$
trên miền nghiệm có $\COV\geq\COV^*$.

Với mỗi mức suy giảm cho phép
\begin{equation}
  \delta\in\{0,0.01,0.025,0.05,0.10\},
\end{equation}
thêm ràng buộc:
\begin{equation}
  \SIM\geq L_\delta
  =\left\lceil(1-\delta)\SIM^*\right\rceil,
  \label{eq:similarity-budget}
\end{equation}
rồi lần lượt tối thiểu hóa $\CONT$, tối thiểu hóa $\OT$ và cuối cùng tối đa hóa
lại $\SIM$ trong các ties. Dùng ceiling bảo đảm tổn thất similarity không vượt
budget đã công bố; dùng floor có thể vi phạm budget gần một đơn vị reward.

Đây là $\varepsilon$-constraint tri-objective với lexicographic completion,
không phải thuật toán liệt kê toàn bộ Pareto front. Sau ba tầng completion,
không thể tồn tại nghiệm khác có $\SIM$ không nhỏ hơn, $\CONT$ không lớn hơn và
$\OT$ không lớn hơn với ít nhất một bất đẳng thức nghiêm; nếu có, nghiệm đó sẽ
mâu thuẫn với một trong ba optimum đã cố định. Do đó mỗi điểm hoàn tất là một
điểm không bị trội, nhưng năm mức $\delta$ chỉ tạo một tập đại diện rời rạc.
Các điểm trùng nhau được gộp; nếu còn budget tính toán, chèn thêm mức $\delta$
vào khoảng có gap lớn nhất trong objective space \cite{mesquita2023}.

\section{Kiến trúc phần mềm đề xuất}

Mã nguồn baseline được giữ làm mốc audit; phần mở rộng C++ nằm trong các tệp
riêng và được link vào executable mới. Cấu trúc triển khai:

\begin{verbatim}
src/
  hcorap2sat.cpp              # executable gốc để audit
  hcorap_multi.cpp            # executable C++ cho B0/B1/B2
  encodings/HCORAP/
    HCORAPEncoding.*          # encoding gốc
  proposed/
    cpp/encodings/            # C++ ablation
      CardinalityNetwork.*
      ImpliedConstraints.*
      SymmetryBreaking.*
      HCORAPMultiObjectiveEncoding.*
    hcorap/                   # Python oracle

experiments/
  run_cpp_experiments.sh
  collect_cpp_results.py
  configs/cpp/
    full_configuration_matrix.env
  results/

docs/
  FAIR_EXPERIMENT_PROTOCOL.md
  IMPLIED_CONSTRAINTS.md
  SYMMETRY_BREAKING.md

instances/
  official/                   # 800 official instances
  corrected_v2/
  stress/

tests/
  instances/tradeoff.txt
  instances/symmetry.txt
  instances/symmetry_partial.txt
  test_cpp_campaign.py
  test_cpp_multiobjective.py
  test_cpp_crosscheck.py
  ...                         # unit tests cho oracle/generator
\end{verbatim}

Phép so sánh phương pháp chính dùng chung parser, hard model, WCNF writer,
compiler flags và Open-WBO C++ backend. Python không nằm trên đường chạy được
đo. Trước full experiment, B0 C++ mới phải khớp optimum của C++ gốc và brute
force/RC2 oracle trên các instance nhỏ; assignment có thể khác khi đồng tối ưu.

Kiểm thử triển khai phát hiện helper \texttt{addPBGEQ} của thư viện C++ kế thừa
ghi vào vector chưa cấp kích thước và truyền nhầm cận. Để không thay đổi baseline
audit, phần mở rộng không sửa helper cũ mà mã hóa cục bộ
$\sum_i q_i x_i\geq K$ thành
$\sum_i q_i(1-x_i)\leq\sum_i q_i-K$. Mọi bound coverage/similarity của B1--B2
đều dùng phép biến đổi này và có regression test.

\subsection{Ba trục cấu hình encoding đã triển khai}

Các cải tiến encoding được thiết kế thành ba trục độc lập. Giá trị đầu tiên ở
mỗi trục là baseline mặc định:

\begin{table}[H]
\centering
\caption{Ma trận cấu hình encoding trong executable \texttt{hcorap\_multi}.}
\begin{tabularx}{\textwidth}{@{}p{3.4cm}>{\raggedright\arraybackslash}p{4.4cm}Y@{}}
\toprule
Trục & Giá trị & Vai trò \\
\midrule
Cardinality
  & \texttt{sorting-network}, \texttt{totalizer}
  & So sánh sorting network gốc với Totalizer hai chiều. \\
Implied constraints
  & \texttt{none}, \texttt{user-slots}, \texttt{slot-capacity},
    \texttt{both}, \texttt{both-plus}
  & Bổ sung các hệ quả logic và projection để tăng propagation. \\
Symmetry breaking
  & \texttt{none}, \texttt{slots}, \texttt{services},
    \texttt{slot-service}, \texttt{all}
  & Chọn một đại diện canonical trong mỗi lớp đối xứng exact. \\
\bottomrule
\end{tabularx}
\end{table}

Totalizer được mã hóa hai chiều để mỗi output unary tương đương chính xác với
``có ít nhất $k$ input đúng''. Vì vậy các biến threshold dùng trong overtime và
các tầng lexicographic không thể nhận giá trị tùy ý. Sorting network vẫn là
baseline mặc định để bảo toàn khả năng audit.

Nhóm implied constraints gồm: (i) đẳng thức giữa số occupied slot của user và
số service đã phục vụ; (ii) upper bound theo maximum matching agent--user tại
mỗi time slot; và (iii) cấu hình \texttt{both-plus} bổ sung projected service
assignment, effective workload cap và service-slot clustering. Trong soft
coverage, user-slot cardinality được channel với tổng biến
\texttt{performed}, không bị cố định bằng tổng số service ban đầu.

Đặt $E(a,s,h)=[r(a,s)>0]\land TSA(a,h)\land TSS(s,h)$ là điều kiện tồn tại
candidate assignment.
Symmetry breaking chỉ gộp các đối tượng tương đương exact. Hai slot phải có
cùng vector candidate $E(a,s,h)$; hai service phải cùng user, cùng sequence,
cùng reward và cùng candidate; hai agent phải cùng $\HN$, $\HE$, reward và
candidate. Value precedence được dùng để chọn thứ tự canonical. Implied
constraints giữ nguyên toàn bộ tập lịch khả thi, trong khi symmetry breaking
loại các lịch có nhãn khác nhau nhưng tương đương và giữ ít nhất một đại diện
trong mỗi orbit; cả hai phải bảo toàn SAT/UNSAT và optimum.

Ma trận đầy đủ có:
\begin{equation}
  2\times5\times5=50\ \text{cấu hình encoding}.
\end{equation}
Runner ghi từng run vào JSON và \path{manifest.tsv}. Sau campaign, script
\path{collect_cpp_results.py} tạo hai bảng UTF-8 BOM mở trực tiếp bằng Excel:
\path{runs.csv} chứa dữ liệu thô và \path{configuration_summary.csv} chứa
solved/timeout/error count, mean, median, PAR-2, encode/solve time, số biến,
hard/soft clauses và quality metrics. Tại thời điểm cập nhật, toàn bộ
112 regression tests đã qua; smoke campaign nhỏ gồm 200 top-level runs trên đủ
50 cấu hình đều đạt \texttt{OPTIMUM} và qua verifier. Đây chỉ là kiểm thử chức
năng bằng RC2-compatible test backend, không phải bằng chứng cải thiện runtime
trên Open-WBO hoặc trên toàn bộ official benchmark.

\section{Các gói công việc}

\subsection{WP1: Đóng băng và tái lập baseline}

\textbf{Mục tiêu:} xác nhận C++ gốc, parser, objective và bộ kiểm chứng nghiệm.

\textbf{Nhiệm vụ:}
\begin{enumerate}
  \item tạo tag hoặc branch \texttt{baseline-original};
  \item ghi lại commit hash, compiler và tham số build;
  \item biên dịch \texttt{hcorap2sat} và sinh WCNF;
  \item chạy một tập đại diện bằng solver được lựa chọn;
  \item xây dựng verifier độc lập để tính $\COV$, $\SIM$, $\CONT$ và $\OT$;
  \item tái lập trạng thái SAT/UNSAT/optimality và quality metrics của toàn bộ
        800 instance; tái tạo đầy đủ các ô của Bảng 5--6 trong phạm vi solver
        và time limit khả dụng.
\end{enumerate}

Để so sánh trực tiếp, báo cáo cả statistic gốc (runtime trung bình trên các
instance được certified) và số liệu đã sửa gồm PAR-2, cactus plot, solved count
và timeout count. Hai loại statistic phải đặt cạnh nhau và ghi rõ rằng khác
biệt do PAR-2 không phải lỗi implementation. Không yêu cầu runtime tuyệt đối
khớp bài báo vì phần cứng và phiên bản solver khác nhau; yêu cầu bắt buộc là
khớp feasibility, objective và xu hướng theo cấu hình khi điều kiện cho phép.

\textbf{Kiểm thử:} tạo các instance rất nhỏ có thể vét cạn toàn bộ nghiệm:
2 agents--2 services; một sequence; một trường hợp cần overtime; một trường
hợp UNSAT; và một trường hợp có ít nhất hai nghiệm không trội.

\textbf{Điều kiện hoàn thành:}
\begin{itemize}
  \item mọi nghiệm đều thỏa hard constraints;
  \item objective do verifier tính khớp objective của solver;
  \item B0 C++ mở rộng và C++ gốc đạt cùng optimum trên các instance nhỏ;
  \item oracle độc lập đạt cùng optimum và verifier chấp nhận mọi nghiệm;
  \item khác biệt về assignment chỉ xuất hiện khi có nhiều nghiệm tối ưu.
\end{itemize}

\subsection{WP2: Xây dựng benchmark hiệu chỉnh}

Hai bộ benchmark được duy trì song song:
\begin{itemize}
  \item \textbf{Official}: giữ nguyên 800 instance để tái lập và so sánh trực tiếp;
  \item \textbf{Corrected v2}: sửa generator và thiết kế paired/nested.
\end{itemize}

Generator v2 phải:
\begin{enumerate}
  \item sinh language từ đúng tập \texttt{languages\_labels};
  \item biểu diễn language nhất quán bằng danh sách;
  \item bảo đảm mỗi user có đúng $V$ services nếu $V$ được gọi là
        services-per-user;
  \item sinh 25 agents trước và lấy các tập lồng nhau
        $A_{10}\subset A_{15}\subset A_{20}\subset A_{25}$;
  \item tạo cấu hình $V=4$ như tập con của cấu hình $V=5$;
  \item sử dụng và lưu random seed;
  \item lưu raw attributes, raw similarity và ma trận reward đã rời rạc hóa;
  \item gắn nhãn nguyên nhân infeasibility, đặc biệt là service không có candidate;
  \item tạo các mức tải low, critical và overload.
\end{enumerate}

Có thể điều khiển mức tải sơ bộ bằng:
\begin{equation}
  \rho=\frac{S}{\sum_a(\HN(a)+\HE(a))}.
\end{equation}
Tuy nhiên, $\rho$ không phải tiêu chuẩn feasibility: kết quả còn phụ thuộc time
windows, qualification, candidate density và xung đột simultaneity. Vì vậy các
giá trị $0.60$, $0.85$ và $1.05$ chỉ được xem là điểm khởi đầu của grid, không
được gắn nhãn low/critical/overload trước khi có dữ liệu.

Cuối tuần 5 phải chạy micro-pilot trên \emph{calibration seeds} không xuất hiện
trong tập đánh giá cuối. Ngoài $\rho$, ghi nhận số candidate agent--time pairs
trên mỗi service, tỷ lệ service có đúng một candidate và phân phối
qualification scarcity. Nhãn tải tạm thời được gán theo hành vi baseline:
\begin{itemize}
  \item low: ít nhất $90\%$ instance full-coverage feasible và dưới $10\%$
        nghiệm tối ưu có overtime;
  \item critical: $60$--$90\%$ instance full-coverage feasible và ít nhất
        $30\%$ nghiệm feasible có overtime;
  \item overload stress test: ít nhất $80\%$ instance infeasible khi coverage
        là hard, nhưng formulation coverage mềm vẫn đạt $\COV_N\geq0.85$.
\end{itemize}
Các ngưỡng này là quy tắc thiết kế provisional; phải được freeze cùng generator
trước khi mở evaluation seeds và không được chỉnh theo kết quả của B1/B2.

Với $U\in\{30,40\}$, $A\in\{10,15,20,25\}$,
$V\in\{4,5\}$, ba vùng tham số và 5 calibration seeds, micro-pilot tối đa gồm:
\begin{equation}
  2\times4\times2\times3\times5=240\ \text{instance evaluations}.
\end{equation}
Micro-pilot chỉ dùng hard-feasibility check và B0 gốc với timeout ngắn; không
chạy toàn bộ 16 weights hoặc mọi tầng multi-objective. Sau khi freeze benchmark,
số evaluation seeds được quyết định bằng precision/power analysis ở WP6.

\subsection{WP3: Weighted sensitivity analysis}

Không chạy cơ học $16\times800=12{,}800$ solver calls trên official benchmark.
Thiết kế tuần tự gồm:
\begin{enumerate}
  \item chạy B0 gốc $(1,1)$ trên toàn bộ 800 official instances;
  \item chạy 16 cặp trọng số trên subset official và corrected-calibration được
        phân tầng theo kích thước, feasibility, candidate density và load;
  \item dùng kết quả screening để chọn tối đa 4--6 cấu hình đại diện hoặc tinh
        chỉnh quanh các weight breakpoints;
  \item chỉ chạy các cấu hình đã freeze trên evaluation set đầy đủ.
\end{enumerate}

Trước mỗi campaign phải tính
\begin{equation}
  N_{\mathrm{call}}=\sum_m N_{\mathrm{inst},m}
                       N_{\mathrm{stage},m}
\end{equation}
và core-hour worst case theo cumulative timeout. Campaign vượt compute budget
phải giảm subset hoặc số cấu hình trước khi chạy, không được bỏ timeout sau khi
đã nhìn kết quả. Phân tích:
\begin{itemize}
  \item thay đổi assignment giữa các cấu hình trọng số;
  \item thay đổi $\SIM_N$, $\CONT_N$ và $\OT_N$;
  \item tỷ lệ instance có nhiều nghiệm chất lượng khác nhau nhưng weighted
        objective gần tương đương;
  \item mức độ trade-off theo load và số agent.
\end{itemize}

\subsection{WP3B: Ablation encoding và symmetry breaking}

\textbf{Mục tiêu:} đo riêng ảnh hưởng của cách mã hóa và strengthening, không
đồng nhất giảm số clause với giảm thời gian giải.

\textbf{Nhiệm vụ:}
\begin{enumerate}
  \item chạy paired pilot cho hai cardinality encodings khi implied constraints
        và symmetry đều cố định ở \texttt{none};
  \item cố định cardinality encoding rồi ablate lần lượt năm cấu hình implied;
  \item ablate năm cấu hình symmetry trên cùng thứ tự instance và cùng solver;
  \item chỉ chạy ma trận $2\times5\times5$ đầy đủ trên subset phân tầng sau khi
        pilot xác nhận compute budget;
  \item so sánh số biến, hard/soft clauses, encode time, solve time, solved count
        và PAR-2; xác nhận quality metrics và optimum không thay đổi;
  \item chọn tối đa một hoặc hai cấu hình encoding cho bảng so sánh policy chính,
        đồng thời luôn giữ baseline \texttt{sorting-network/none/none}.
\end{enumerate}

Symmetry breaking được đánh giá riêng vì nó chủ động loại các assignment có
nhãn khác nhau trong cùng orbit. Regression phải chứng minh còn ít nhất một đại
diện cho mỗi orbit trên tiny instances, giữ nguyên optimum trong cả full và soft
coverage, và không tạo overhead encoding khi detector không tìm thấy lớp đối xứng.

\subsection{WP4: Triển khai Lexicographic MaxSAT}

Phiên bản đầu có thể khởi tạo lại solver giữa các tầng để ưu tiên tính đúng.
Phiên bản sau tái sử dụng solver state bằng assumptions/incremental constraints
nếu thư viện hỗ trợ ổn định.

Timeout được áp dụng cho \emph{toàn bộ policy trên một instance}, không nhân
riêng cho từng tầng. Full-coverage instances cần ba tầng; overload stress test
cần thêm tầng coverage. Runtime có thể nhỏ hơn hoặc lớn hơn bội số số tầng vì
các bounds làm thay đổi search space, nên phải báo cáo cumulative time, time mỗi
tầng và số solver calls thay vì giả định overhead đúng $3\times$ hoặc $4\times$.
Go/no-go gate cuối tuần 9: nếu dưới $60\%$ pilot instances hoàn tất trong
cumulative timeout 3600 giây, B1 chỉ chạy trên critical subset hoặc chuyển sang
incremental/native multi-objective implementation trước full experiment.

Kiểm thử bắt buộc:
\begin{itemize}
  \item so sánh kết quả với brute force;
  \item chứng minh tầng sau không làm xấu giá trị tối ưu của tầng trước;
  \item kiểm tra cả hai thứ tự continuity-priority và overtime-priority;
  \item kiểm tra objective có optimum bằng 0;
  \item kiểm tra nhiều nghiệm đồng tối ưu ở mỗi tầng;
  \item kiểm tra instance infeasible khi coverage là hard constraint.
\end{itemize}

\subsection{WP5: Triển khai similarity-budget
\texorpdfstring{$\varepsilon$}{epsilon}-constraint}

Với mỗi instance:
\begin{enumerate}
  \item tính $\SIM^*$;
  \item lần lượt thêm năm mức similarity budget trong
        Eq.~\eqref{eq:similarity-budget};
  \item tối thiểu hóa $\CONT$, thêm $\CONT\leq\CONT^*$; tối thiểu hóa $\OT$,
        thêm $\OT\leq\OT^*$; cuối cùng tối đa hóa lại $\SIM$;
  \item lưu toàn bộ nghiệm và metric;
  \item kiểm chứng non-dominance, gộp điểm trùng và ghi realized similarity loss.
\end{enumerate}

\subsection{WP6: Thực nghiệm đầy đủ và phân tích thống kê}

Chỉ chạy thực nghiệm đầy đủ sau khi pilot xác nhận benchmark có đủ ba mức độ
khó và overtime thực sự được kích hoạt trên critical/overload instances.
Calibration seeds chỉ ước lượng variance và tỷ lệ zero differences. Trước khi
mở evaluation seeds, xác định smallest effect size of interest (SESOI), mô phỏng
power của paired Wilcoxon/sign test dưới các cỡ mẫu ứng viên và chọn số seed để
đạt power mục tiêu ít nhất $0.8$ nếu compute budget cho phép. Nếu không đạt,
nghiên cứu phải báo confidence interval/độ chính xác ước lượng và không dùng
``không có ý nghĩa thống kê'' như bằng chứng hai phương pháp tương đương.

\section{Thiết kế thực nghiệm}

\subsection{Các phương pháp so sánh}

\begin{table}[H]
\centering
\caption{Các phương pháp trong thực nghiệm.}
\begin{tabularx}{\textwidth}{@{}>{\raggedright\arraybackslash}p{3.4cm}YY@{}}
\toprule
Phương pháp & Mục đích & Ghi chú \\
\midrule
C++ gốc & Audit/tái lập baseline bài báo & Báo ở bảng tái lập riêng \\
C++ weighted B0 & B0 gốc và calibrated weighted & Cùng model/writer/backend với B1--B2 \\
C++ Lex (continuity) & Chính sách ưu tiên continuity & Sequential MaxSAT C++ \\
C++ Lex (overtime) & Chính sách ưu tiên overtime & Sequential MaxSAT C++ \\
C++ $\varepsilon$-constraint & Trade-off có similarity budget & Năm mức chính, có thể refine \\
C++ CP-SAT hoặc MILP & Baseline ngoài MaxSAT & Chỉ đưa vào runtime chính sau khi port C++ \\
\bottomrule
\end{tabularx}
\end{table}

\subsection{Ma trận cấu hình thực nghiệm}

Ba trục encoding độc lập với trục policy. File
\path{full_configuration_matrix.env} khai báo 50 tổ hợp; runner duyệt tuần
tự cardinality--implied--symmetry và, trong mỗi tổ hợp, chạy weighted,
lex-continuity, lex-overtime cùng các mức $\delta$ của
$\varepsilon$-constraint. Với năm mức $\delta$, một instance cần tối đa:
\begin{equation}
  50\times(3+5)=400\ \text{top-level runs},
\end{equation}
chưa tính các solver stage nội bộ của từng policy. Vì vậy full matrix chỉ dùng
cho screening subset; bảng policy chính phải dùng encoding đã freeze từ pilot.

Mỗi JSON lưu \path{cardinality_encoding},
\path{implied_constraints}, \path{symmetry_breaking}, method, delta,
timing theo stage, số biến/clause và metrics đã verify. Manifest lưu thêm
SHA-256 của instance, result path và exit code. Không dùng lại result directory
khi schema hoặc campaign thay đổi.

\subsection{Môi trường chạy}

Mọi phương pháp phải dùng:
\begin{itemize}
  \item cùng ngôn ngữ C++ trên đường chạy được đo và cùng mức tối ưu compiler;
  \item cùng parser, hard model, WCNF writer và MaxSAT backend cho B0--B2;
  \item cùng một CPU core và cùng máy;
  \item cùng giới hạn RAM;
  \item cùng timeout;
  \item cùng instance và cùng objective definition;
  \item ít nhất ba lần chạy nếu solver không deterministic;
  \item log đầy đủ phiên bản compiler, solver, OS và tham số dòng lệnh.
\end{itemize}

Thời gian end-to-end được tính từ lúc đọc instance đến sau khi verifier C++
kiểm tra nghiệm, gồm parse, encode, serialize WCNF, solve và verify. Timeout là
ngân sách tích lũy cho toàn bộ policy. Shell runner chỉ điều phối và không chứa
implementation của phương pháp. Tái lập runtime của executable C++ gốc được
báo riêng nếu encoding hoặc định dạng WCNF chưa hoàn toàn đồng nhất.

Các mốc timeout đề xuất là 60, 300 và 3600 giây.

\subsection{KPI tính toán}

\begin{itemize}
  \item số instance tìm được nghiệm feasible;
  \item số instance chứng nhận optimal/UNSAT;
  \item time-to-first-solution;
  \item time-to-best-solution;
  \item time-to-proof;
  \item peak memory;
  \item PAR-2;
  \item số biến và số clauses của encoding.
\end{itemize}

\subsection{KPI chất lượng nghiệm}

\begin{itemize}
  \item $\COV_N$, $\SIM_N$, $\CONT_N$ và $\OT_N$;
  \item số caregiver trung bình và cực đại trên mỗi sequence;
  \item số agent có overtime;
  \item workload trung bình, độ lệch chuẩn và cực đại;
  \item số điểm đại diện không bị trội và số budget tạo nghiệm trùng;
  \item mức cải thiện continuity/overtime tại mỗi similarity budget.
\end{itemize}

\subsection{Phương pháp thống kê}

Do corrected benchmark được thiết kế paired/nested, sử dụng:
\begin{itemize}
  \item median và interquartile range thay cho chỉ báo average đơn lẻ;
  \item bootstrap $95\%$ confidence interval;
  \item paired Wilcoxon signed-rank test;
  \item paired rank-biserial effect size và smallest effect size of interest;
  \item pilot-based power analysis và Holm correction cho multiple comparisons;
  \item cactus plot và solved-over-time curve cho runtime;
  \item objective-space trade-off plot cho ba thành phần chất lượng nghiệm.
\end{itemize}

Đơn vị thống kê độc lập là base scenario/seed. Các biến thể nested theo $A$ và
$V$ trong cùng seed là repeated measures, không được gộp như những quan sát độc
lập. Nếu paired differences có quá nhiều giá trị 0, dùng Wilcoxon--Pratt hoặc
paired sign test và báo rõ quy ước. P-value luôn đi cùng effect size và
confidence interval.

Không tính runtime trung bình chỉ trên các instance đã giải xong, vì cách này
loại bỏ các timeout và gây survivor bias.

\section{Ablation study}

Các cấu hình tối thiểu cần chạy:
\begin{enumerate}
  \item weighted objective gốc trên official benchmark;
  \item weighted objective gốc với metric chuẩn hóa;
  \item weighted sensitivity trên official benchmark;
  \item weighted sensitivity trên corrected benchmark;
  \item lexicographic trên official benchmark;
  \item lexicographic trên corrected benchmark;
  \item $\varepsilon$-constraint trên critical-load subset;
  \item overload stress test với coverage là mục tiêu ưu tiên cao nhất.
  \item sorting network so với Totalizer khi các trục khác ở baseline;
  \item năm mức implied constraints với cardinality và symmetry cố định;
  \item năm mức symmetry breaking với cardinality và implied cố định;
  \item ma trận $2\times5\times5$ trên screening subset để phát hiện interaction.
\end{enumerate}

Ablation giúp phân biệt cải tiến đến từ phương pháp tối ưu, benchmark mới hay
cách chuẩn hóa metric. Ablation encoding còn phân biệt ảnh hưởng của cấu trúc
CNF, redundant constraints và loại bỏ nghiệm đối xứng. Không chọn cấu hình chỉ
dựa trên một instance hoặc chỉ dựa trên số clause.

\section{Tiêu chí thành công}

\subsection{Mức tối thiểu}

\begin{itemize}
  \item tái lập được feasibility/optimality của baseline trên một tập đại diện;
  \item verifier độc lập khớp objective trên toàn bộ nghiệm đã chạy;
  \item chứng minh similarity tổng và similarity/service dẫn tới diễn giải khác nhau;
  \item định lượng được độ nhạy của weighted solution theo trọng số;
  \item sinh được ít nhất hai điểm không bị trội trên một tỷ lệ đáng kể các
        critical-load instances.
\end{itemize}

\subsection{Mức tốt}

Mục tiêu thực nghiệm, không phải kết quả được giả định trước:
\begin{itemize}
  \item tại similarity loss không quá $2.5\%$, confidence interval của thay đổi
        continuity được báo cùng SESOI; chỉ tuyên bố cải thiện khi cả độ lớn
        thực tiễn và độ bất định đều đạt tiêu chí đã preregister;
  \item Lexicographic/$\varepsilon$-constraint MaxSAT giải được phần lớn pilot instances trong
        300--3600 giây;
  \item các kết luận chính ổn định qua nhiều seeds và có confidence interval hẹp;
  \item benchmark v2 tạo được đủ instance low, critical và overload theo thiết kế.
\end{itemize}

\subsection{Mức đủ mạnh để công bố}

\begin{itemize}
  \item benchmark v2 và generator có thể tái lập;
  \item implementation C++ weighted/lexicographic/$\varepsilon$-constraint
        và verifier độc lập;
  \item baseline MaxSAT và ít nhất một baseline ngoài MaxSAT;
  \item tập điểm trade-off không bị trội, có phân tích paired và effect size;
  \item source code, dữ liệu, cấu hình và raw result được công bố đầy đủ.
\end{itemize}

\section{Rủi ro và phương án xử lý}

\begin{longtable}{@{}p{0.31\textwidth}p{0.62\textwidth}@{}}
\caption{Rủi ro nghiên cứu và phương án xử lý.}\\
\toprule
Rủi ro & Phương án xử lý \\
\midrule
\endfirsthead
\toprule
Rủi ro & Phương án xử lý \\
\midrule
\endhead
Official benchmark gần như không có overtime & Tạo critical-load và overload
suite; chỉ dùng official benchmark cho tái lập. \\
$\varepsilon$-constraint quá chậm & Áp cumulative timeout; giới hạn năm budget
chính; tái sử dụng solver state và chỉ refine khi còn compute budget. \\
C++ B0 mở rộng không khớp C++ gốc & Dừng campaign; dùng brute force, RC2 oracle
và verifier độc lập để định vị khác biệt objective/encoding. \\
Corrected benchmark quá dễ hoặc quá khó & Chạy pilot, hiệu chỉnh load theo
feasibility thực tế thay vì chỉ dùng tỷ lệ capacity. \\
Continuity cải thiện rất ít & Tăng tỷ lệ sequence nhiều services và tạo
qualification scarcity có kiểm soát. \\
EvalMaxSAT không ổn định trên nền tảng chạy & Khóa Open-WBO C++ đã qua smoke
test cho mọi phương pháp; không thay solver giữa campaign. \\
Symmetry breaking loại nhầm nghiệm không tương đương & Chỉ dùng signature exact
trên reward, candidate, \texttt{SU/SEQ}, $\HN/\HE$; kiểm tra optimum bằng RC2/brute force
trên full và soft coverage trước campaign. \\
Ma trận $2\times5\times5$ làm vượt compute budget & Chạy tuần tự các ablation
một trục, screening trên subset phân tầng, rồi mới freeze một số cấu hình cho
evaluation set. \\
Không có dữ liệu vận hành thật & Giới hạn kết luận ở benchmark mô phỏng; không
tuyên bố hiệu quả lâm sàng hoặc triển khai thực tế. \\
Reviewer xem sửa generator là đóng góp nhỏ & Xác định benchmark là điều kiện
nền; đóng góp chính là nghiên cứu tách biệt benchmark/weights/optimizer, so sánh
với calibrated weighted và baseline ngoài MaxSAT, cùng phân tích trade-off. \\
Calibrated weighted đã đủ tốt & Báo đây là kết quả âm có giá trị: hạn chế nằm ở
benchmark/đánh giá; không thổi phồng lexicographic như một cải tiến mặc định. \\
Thiếu power hoặc compute budget & Giảm số policy bằng screening, tăng seed trước
khi tăng số weight, báo precision/CI và không kết luận tương đương từ $p>0.05$. \\
\bottomrule
\end{longtable}

\clearpage
\section{Tiến độ 14 tuần và 2 tuần dự phòng}

\begin{table}[H]
\centering
\caption{Kế hoạch triển khai có stage gates.}
\begin{tabularx}{\textwidth}{@{}p{2cm}Yp{4cm}@{}}
\toprule
Thời gian & Công việc & Sản phẩm \\
\midrule
Tuần 1--2 & Đóng băng C++ gốc, viết verifier và brute-force tests
           & Baseline và bộ test nhỏ \\
Tuần 3 & Tái lập Bảng 5--6; báo cả statistic gốc và PAR-2
           & Báo cáo tái lập \\
Tuần 4 & Sửa generator, thiết kế nested benchmark và diagnostics
           & Generator v2 alpha \\
Tuần 5 & Chạy micro-pilot trên calibration seeds; hiệu chỉnh và freeze generator
           & Gate G1: benchmark v2 frozen \\
Tuần 6 & Paired pilot Totalizer, implied constraints và symmetry breaking
           & Encoding shortlist; ma trận/log đã kiểm chứng \\
Tuần 7 & Weighted screening, adaptive refinement và chốt compute budget
           & Calibrated B0; policy shortlist \\
Tuần 8--9 & Triển khai, kiểm thử và benchmark sequential Lexicographic MaxSAT
           & Gate G2: B1 scalability \\
Tuần 10 & Triển khai $\varepsilon$-constraint và kiểm chứng non-dominance
           & B2 solver \\
Tuần 11 & Integration pilot; power/precision analysis; freeze protocol
           & Gate G3: preregistered protocol \\
Tuần 12--13 & Chạy thực nghiệm đầy đủ và rerun lỗi kỹ thuật
           & Raw results, logs, verified solutions \\
Tuần 14 & Phân tích thống kê sơ bộ và audit kết quả
           & Tables, figures, audit report \\
Tuần 15--16 & Contingency cho benchmark/solver; nếu không dùng thì viết bài và
               đóng gói reproducibility package
           & Bản thảo và release package \\
\bottomrule
\end{tabularx}
\end{table}

\section{Sản phẩm bàn giao}

Nghiên cứu dự kiến tạo ra:
\begin{enumerate}
  \item bản tái lập baseline C++ gốc;
  \item báo cáo kiểm toán official benchmark;
  \item corrected HCORAP benchmark v2 và generator có seed;
  \item weighted, lexicographic và $\varepsilon$-constraint MaxSAT
        implementations bằng C++;
  \item hai cardinality encodings cùng năm cấu hình implied constraints và năm
        cấu hình symmetry breaking có thể ablate độc lập;
  \item independent solution verifier;
  \item unit tests và tiny brute-force instances;
  \item ma trận cấu hình, manifest, raw CSV tương thích Excel, summary PAR-2,
        logs và script tạo biểu đồ;
  \item cactus plots, normalized-quality plots và objective-space trade-off plots;
  \item bản thảo bài báo hoặc luận văn;
  \item reproducibility package kèm hướng dẫn chạy.
\end{enumerate}

\section{Thông điệp khoa học dự kiến}

Thông điệp trung tâm của nghiên cứu được phát biểu thận trọng như sau:

\begin{quote}
Weighted MaxSAT cung cấp cách biểu diễn tự nhiên cho HCORAP, nhưng kết luận về
chất lượng của nó phải tách ảnh hưởng của benchmark và cách hiệu chỉnh trọng số.
Trên benchmark được kiểm soát, calibrated weighted, sequential lexicographic và
$\varepsilon$-constraint MaxSAT cho các semantics quyết định khác nhau. Nghiên
cứu xác định khi nào sự khác biệt đó có ý nghĩa thực tiễn, mức chi phí tính toán
đi kèm, và chấp nhận kết quả rằng calibrated weighted đã đủ tốt nếu dữ liệu ủng
hộ kết luận này.
\end{quote}

Kết luận cuối cùng chỉ được chấp nhận nếu được hỗ trợ bởi thực nghiệm paired,
metric chuẩn hóa, kiểm định thống kê và bộ kiểm chứng nghiệm độc lập.

\section{Hướng mở rộng sau nghiên cứu}

Sau khi hoàn thành hướng chính, có thể mở rộng theo thứ tự:
\begin{enumerate}
  \item adaptive selection giữa cardinality/implied/symmetry configurations
        dựa trên đặc trưng instance, cùng các AMO encodings mới;
  \item UNSAT-core explanation và minimal resource augmentation;
  \item incremental MaxSAT cho dynamic rescheduling và minimum disruption;
  \item phối hợp specialist giữa nhiều khu vực;
  \item tích hợp travel time, service duration và routing;
  \item kiểm thử trên dữ liệu vận hành thật.
\end{enumerate}

\begin{thebibliography}{9}

\bibitem{unceta2024}
I. Unceta, B. Salbanya, J. Coll, M. Villaret, and J. Nin,
``Optimizing resource allocation in home care services using MaxSAT,''
\emph{Cognitive Systems Research}, vol. 88, 101291, 2024.
\href{https://doi.org/10.1016/j.cogsys.2024.101291}
{doi:10.1016/j.cogsys.2024.101291}.

\bibitem{salbanya2023}
B. Salbanya, I. Unceta, and J. Nin,
``Resource Allocation in Home Care Services Using Reinforcement Learning,''
in \emph{Artificial Intelligence Research and Development}, 2023.
\href{https://doi.org/10.3233/FAIA230680}{doi:10.3233/FAIA230680}.

\bibitem{greenhhc2024}
``A multiobjective approach for weekly Green Home Health Care routing and
scheduling problem with care continuity and synchronized services,''
\emph{Operations Research Perspectives}, vol. 12, 100302, 2024.
\href{https://doi.org/10.1016/j.orp.2024.100302}
{doi:10.1016/j.orp.2024.100302}.

\bibitem{reoptimization2024}
``An exact two-phase approach to re-optimize tours in home care planning,''
\emph{Computers \& Operations Research}, vol. 161, 106408, 2024.
\href{https://doi.org/10.1016/j.cor.2023.106408}
{doi:10.1016/j.cor.2023.106408}.

\bibitem{demirovic2019}
E. Demirovi\'c, N. Musliu, and F. Winter,
``Modeling and solving staff scheduling with partial weighted MaxSAT,''
\emph{Annals of Operations Research}, vol. 275, pp. 79--99, 2019.
\href{https://doi.org/10.1007/s10479-017-2693-y}
{doi:10.1007/s10479-017-2693-y}.

\bibitem{mosquera2019}
F. Mosquera, P. Smet, and G. Vanden Berghe,
``Flexible home care scheduling,''
\emph{Omega}, vol. 83, pp. 80--95, 2019.
\href{https://doi.org/10.1016/j.omega.2018.02.005}
{doi:10.1016/j.omega.2018.02.005}.

\bibitem{bodvarsdottir2025}
E. B. B\"o\dh varsd\'ottir and T. Stidsen,
``A review of multi-objective optimization methods for personnel rostering
problems,'' \emph{Journal of Scheduling}, 2025.
\href{https://doi.org/10.1007/s10951-025-00845-0}
{doi:10.1007/s10951-025-00845-0}.

\bibitem{mesquita2023}
M. Mesquita-Cunha, J. R. Figueira, and A. P. Barbosa-P\'ovoa,
``New $\varepsilon$-constraint methods for multi-objective integer linear
programming: A Pareto front representation approach,''
\emph{European Journal of Operational Research}, vol. 306, no. 1,
pp. 286--307, 2023.
\href{https://doi.org/10.1016/j.ejor.2022.07.044}
{doi:10.1016/j.ejor.2022.07.044}.

\bibitem{jabs2025}
C. Jabs, J. Berg, B. Bogaerts, and M. J\"arvisalo,
``Certifying Pareto-optimality in multi-objective maximum satisfiability,''
2025, arXiv:2501.17493.
\href{https://doi.org/10.48550/arXiv.2501.17493}
{doi:10.48550/arXiv.2501.17493}.

\end{thebibliography}

\end{document}
